\section{Training}

The training pipeline of LongCat-Video-Avatar 1.5 consists of three 
progressive stages: Base Model Training, RLHF Training, 
and Acceleration Training.  
The first stage establishes the foundational capability of audio-driven 
avatar generation. The model is trained to synthesize temporally coherent 
and identity-preserving video conditioned on speech signals. 
The second stage leverages Reinforcement Learning from Human Feedback 
(RLHF)~\cite{ouyang2022training}, specifically employing Group Relative 
Policy Optimization (GRPO)~\cite{shao2024deepseekmath} to align the 
model's outputs with human preferences. This stage improves perceptual 
quality, expressiveness, and overall generation fidelity beyond what 
supervised training alone can achieve, ensuring that synthesized avatars 
are more natural and visually appealing to human observers.
The third stage focuses on optimizing inference efficiency. Through 
dedicated acceleration training, the model achieves high-quality avatar 
video synthesis at significantly reduced computational cost, enabling 
practical deployment without sacrificing generation quality.

\noindent\textbf{Base Model Training.}
We adopt the flow matching framework~\cite{lipman2022flow} for the 
generative process. Given a clean video latent $x_0$, a noise 
sample $\epsilon \sim \mathcal{N}(\mathbf{0}, \mathbf{I})$, 
and a timestep $t \in [0, 1]$, the noisy latent $x_t$ is 
constructed via linear interpolation:
\begin{equation}
    x_t = (1 - t) \cdot x_0 + t \cdot \epsilon.
    \label{eq:xt}
\end{equation}
The network is trained to predict the velocity 
$v_{\text{pred}}(x_t, c, t; \theta)$, 
where $c$ denotes the task conditions (\emph{e.g.}, text prompts, audios  
and conditional image/video latents), and $\theta$ denotes
the model parameters. The training objective minimizes MSE against
 ground truth velocity $v_t = x_0 - \epsilon$:
\begin{equation}
    \mathcal{L} = \mathbb{E}_{\epsilon, x_0, c, t}
    \left\| v_{\text{pred}}(x_t, c, t; 
    \theta) - v_t \right\|^2.
    \label{eq:loss}
\end{equation}
The base model training consists of multiple progressive stages,  as
outlined in Table~\ref{tab:training_stages}. the training begins with low-resolution pretraining, where the model 
learns the fundamental correspondence between speech signals and facial 
dynamics at a coarse spatial scale, establishing the core audio-driven 
generation capability.
Once the model demonstrates stable audio-driven capability, the training 
transitions to a high-resolution stage, enabling the model to synthesize 
fine-grained visual details and produce high-fidelity avatar videos with 
improved spatial quality.
A reference image module is subsequently introduced into the training 
pipeline, allowing the model to incorporate identity and appearance 
information from a given reference image. This stage equips the model 
with the ability to generate identity-preserving avatars conditioned on 
arbitrary reference inputs.
Finally, the model is trained on a large-scale multi-person dialogue 
dataset, extending its capability to handle multi-person conversational 
scenarios where multiple identities interact in a coherent and temporally 
consistent manner.





\noindent\textbf{RLHF Training.}
Following the base model training, we further improve the model 
performance through a post-training stage. Specifically, we adopt 
the GRPO method described in Section~\ref{sec:grpo}, incorporating 
multiple video quality-related reward signals to guide the optimization process.
We adopt most training hyperparameters and optimization settings from LongCat-Video~\cite{team2025longcat}. For the proposed multi-clip extension, we set the maximum rollout length to 5 clips and randomly sample the actual number of sequential clips during training. To keep training tractable, only the final clip contributes to GRPO optimization, while earlier clips are used solely as temporal context for subsequent rollout. For image-to-video, and video-continuation tasks, we further perform a first-frame hand-presence check with MediaPipe hand detection, which increases the proportion of hand-relevant training samples.




\noindent\textbf{Acceleration Training.}
We adopt the method described in Section~\ref{sec:fsg} for model distillation, 
enabling the distilled model to achieve generation quality comparable 
to 50-step inference using only 8 steps. During this stage, we observe 
that prolonged training leads to a noticeable degradation in visual 
realism. Therefore, we train the model for 400 steps, at which 
point the model achieves the best generation fidelity before 
quality decline sets in. Specifically, the generator learning rate is 
set to $2 \times 10^{-5}$, the fake score learning rate is set to 
$4 \times 10^{-6}$, and the update ratio between the generator and 
the fake scorer is set to $1\!:\!5$.
